% Options for packages loaded elsewhere
\PassOptionsToPackage{unicode}{hyperref}
\PassOptionsToPackage{hyphens}{url}
%
\documentclass[
  12pt,
  a4paper,
  DIV=11,
  numbers=noendperiod]{scrreprt}

\usepackage{amsmath,amssymb}
\usepackage{setspace}
\usepackage{iftex}
\ifPDFTeX
  \usepackage[T1]{fontenc}
  \usepackage[utf8]{inputenc}
  \usepackage{textcomp} % provide euro and other symbols
\else % if luatex or xetex
  \usepackage{unicode-math}
  \defaultfontfeatures{Scale=MatchLowercase}
  \defaultfontfeatures[\rmfamily]{Ligatures=TeX,Scale=1}
\fi
\usepackage{lmodern}
\ifPDFTeX\else  
    % xetex/luatex font selection
\fi
% Use upquote if available, for straight quotes in verbatim environments
\IfFileExists{upquote.sty}{\usepackage{upquote}}{}
\IfFileExists{microtype.sty}{% use microtype if available
  \usepackage[]{microtype}
  \UseMicrotypeSet[protrusion]{basicmath} % disable protrusion for tt fonts
}{}
\usepackage{xcolor}
\setlength{\emergencystretch}{3em} % prevent overfull lines
\setcounter{secnumdepth}{5}
% Make \paragraph and \subparagraph free-standing
\makeatletter
\ifx\paragraph\undefined\else
  \let\oldparagraph\paragraph
  \renewcommand{\paragraph}{
    \@ifstar
      \xxxParagraphStar
      \xxxParagraphNoStar
  }
  \newcommand{\xxxParagraphStar}[1]{\oldparagraph*{#1}\mbox{}}
  \newcommand{\xxxParagraphNoStar}[1]{\oldparagraph{#1}\mbox{}}
\fi
\ifx\subparagraph\undefined\else
  \let\oldsubparagraph\subparagraph
  \renewcommand{\subparagraph}{
    \@ifstar
      \xxxSubParagraphStar
      \xxxSubParagraphNoStar
  }
  \newcommand{\xxxSubParagraphStar}[1]{\oldsubparagraph*{#1}\mbox{}}
  \newcommand{\xxxSubParagraphNoStar}[1]{\oldsubparagraph{#1}\mbox{}}
\fi
\makeatother


\providecommand{\tightlist}{%
  \setlength{\itemsep}{0pt}\setlength{\parskip}{0pt}}\usepackage{longtable,booktabs,array}
\usepackage{calc} % for calculating minipage widths
% Correct order of tables after \paragraph or \subparagraph
\usepackage{etoolbox}
\makeatletter
\patchcmd\longtable{\par}{\if@noskipsec\mbox{}\fi\par}{}{}
\makeatother
% Allow footnotes in longtable head/foot
\IfFileExists{footnotehyper.sty}{\usepackage{footnotehyper}}{\usepackage{footnote}}
\makesavenoteenv{longtable}
\usepackage{graphicx}
\makeatletter
\newsavebox\pandoc@box
\newcommand*\pandocbounded[1]{% scales image to fit in text height/width
  \sbox\pandoc@box{#1}%
  \Gscale@div\@tempa{\textheight}{\dimexpr\ht\pandoc@box+\dp\pandoc@box\relax}%
  \Gscale@div\@tempb{\linewidth}{\wd\pandoc@box}%
  \ifdim\@tempb\p@<\@tempa\p@\let\@tempa\@tempb\fi% select the smaller of both
  \ifdim\@tempa\p@<\p@\scalebox{\@tempa}{\usebox\pandoc@box}%
  \else\usebox{\pandoc@box}%
  \fi%
}
% Set default figure placement to htbp
\def\fps@figure{htbp}
\makeatother

\usepackage{indentfirst}
\usepackage{float}
\floatplacement{figure}{H}
\IfFileExists{plex-otf.sty}{
  %% Full TeXlive
  \usepackage[
    % math,
    RM={Scale=0.94},SS={Scale=0.94},SScon={Scale=0.94},TT={Scale=MatchLowercase,FakeStretch=0.9},
    DefaultFeatures={Ligatures=Common}
  ]{plex-otf}
  % \usepackage[Scale=MatchUppercase]{juliamono}
}{
  %% TinyTeX
  \usepackage{libertine}
}
\KOMAoption{captions}{tableheading}
\makeatletter
\@ifpackageloaded{caption}{}{\usepackage{caption}}
\AtBeginDocument{%
\ifdefined\contentsname
  \renewcommand*\contentsname{Содержание}
\else
  \newcommand\contentsname{Содержание}
\fi
\ifdefined\listfigurename
  \renewcommand*\listfigurename{Список иллюстраций}
\else
  \newcommand\listfigurename{Список иллюстраций}
\fi
\ifdefined\listtablename
  \renewcommand*\listtablename{Список таблиц}
\else
  \newcommand\listtablename{Список таблиц}
\fi
\ifdefined\figurename
  \renewcommand*\figurename{Рисунок}
\else
  \newcommand\figurename{Рисунок}
\fi
\ifdefined\tablename
  \renewcommand*\tablename{Таблица}
\else
  \newcommand\tablename{Таблица}
\fi
}
\@ifpackageloaded{float}{}{\usepackage{float}}
\floatstyle{ruled}
\@ifundefined{c@chapter}{\newfloat{codelisting}{h}{lop}}{\newfloat{codelisting}{h}{lop}[chapter]}
\floatname{codelisting}{Список}
\newcommand*\listoflistings{\listof{codelisting}{Листинги}}
\makeatother
\makeatletter
\makeatother
\makeatletter
\@ifpackageloaded{caption}{}{\usepackage{caption}}
\@ifpackageloaded{subcaption}{}{\usepackage{subcaption}}
\makeatother

\ifLuaTeX
\usepackage[bidi=basic]{babel}
\else
\usepackage[bidi=default]{babel}
\fi
\babelprovide[main,import]{russian}
\babelprovide[import]{english}
% get rid of language-specific shorthands (see #6817):
\let\LanguageShortHands\languageshorthands
\def\languageshorthands#1{}
\usepackage[style=gost-numeric,backend=biber,langhook=extras,autolang=other*]{biblatex}
\addbibresource{bib/cite.bib}
\usepackage{csquotes}
\usepackage{bookmark}

\IfFileExists{xurl.sty}{\usepackage{xurl}}{} % add URL line breaks if available
\urlstyle{same} % disable monospaced font for URLs
\hypersetup{
  pdftitle={Отчёт по лабораторной работе №2},
  pdfauthor={Пакавира Арсениу Висенте Луиш},
  pdflang={ru-RU},
  hidelinks,
  pdfcreator={LaTeX via pandoc}}


\title{Отчёт по лабораторной работе №2}
\usepackage{etoolbox}
\makeatletter
\providecommand{\subtitle}[1]{% add subtitle to \maketitle
  \apptocmd{\@title}{\par {\large #1 \par}}{}{}
}
\makeatother
\subtitle{Дисциплина: Администрирование сетевых подсистем}
\author{Пакавира Арсениу Висенте Луиш}
\date{}

\begin{document}
\maketitle

\renewcommand*\contentsname{Содержание}
{
\setcounter{tocdepth}{1}
\tableofcontents
}
\listoffigures
\listoftables

\setstretch{1.5}
\chapter{Цель
работы}\label{ux446ux435ux43bux44c-ux440ux430ux431ux43eux442ux44b}

Здесь приводится формулировка цели лабораторной работы. Формулировки
цели для каждой лабораторной работы приведены в методических указаниях.

Цель данного шаблона --- максимально упростить подготовку отчётов по
лабораторным работам. Модифицируя данный шаблон, студенты смогут без
труда подготовить отчёт по лабораторным работам, а также познакомиться с
основными возможностями разметки Markdown.

\chapter{Выполнение лабораторной
работы}\label{ux432ux44bux43fux43eux43bux43dux435ux43dux438ux435-ux43bux430ux431ux43eux440ux430ux442ux43eux440ux43dux43eux439-ux440ux430ux431ux43eux442ux44b}

\section{Установка
DNS-сервера}\label{ux443ux441ux442ux430ux43dux43eux432ux43aux430-dns-ux441ux435ux440ux432ux435ux440ux430}

На виртуальной машине server выполнен вход под пользователем user и
переход в режим суперпользователя с помощью команды \texttt{sudo\ -i},
что подтверждается сменой приглашения командной строки на root
(\textbf{?@fig-1}).

\includegraphics[width=0.7\linewidth,height=\textheight,keepaspectratio]{image/1.png}
\href{https://github.com/ArsenioPacavira99/study_2025-2026_net-os-admin-12/tree/master/labs/lab02/report/image/1.png}{рис.
1}

В режиме суперпользователя установлены пакеты \texttt{bind} и
\texttt{bind-utils} с использованием
\texttt{dnf\ -y\ install\ bind\ bind-utils}; транзакция завершена
успешно, зависимости разрешены, пакеты установлены (\textbf{?@fig-2}).

\includegraphics[width=0.7\linewidth,height=\textheight,keepaspectratio]{image/2.png}
\href{https://github.com/ArsenioPacavira99/study_2025-2026_net-os-admin-12/tree/master/labs/lab02/report/image/2.png}{рис.
2}

С помощью утилиты \texttt{dig} выполнен DNS-запрос к имени
\href{http://www.yandex.ru}{www.yandex.ru}; в выводе отображены
заголовок ответа (opcode: QUERY, status: NOERROR), флаги (qr, rd, ra),
секции QUESTION и ANSWER с тремя A-записями и временем отклика, что
подтверждает корректное разрешение имени через внешний DNS-сервер
(\textbf{?@fig-3}).

\includegraphics[width=0.7\linewidth,height=\textheight,keepaspectratio]{image/3.png}
\href{https://github.com/ArsenioPacavira99/study_2025-2026_net-os-admin-12/tree/master/labs/lab02/report/image/3.png}{Нажмите
здесь, чтобы открыть ссылку.} \#\# Конфигурирование кэширующего
DNS-сервера

Проанализировано содержимое файла \texttt{/etc/resolv.conf}: указан
поисковый домен user.net и заданы внешние DNS-серверы (172.249.0.7,
8.8.8.8, 8.8.4.4), что определяет порядок разрешения имён до изменения
конфигурации (\textbf{?@fig-4}).

\includegraphics[width=0.7\linewidth,height=\textheight,keepaspectratio]{image/4.png}
\href{https://github.com/ArsenioPacavira99/study_2025-2026_net-os-admin-12/tree/master/labs/lab02/report/image/4.png}{рис.
4}

Проанализирован файл \texttt{/etc/named.conf}: сервер настроен как
кэширующий (recursion yes), прослушивает порт 53 только на 127.0.0.1,
разрешены запросы от localhost, определены параметры каталогов и
подключены стандартные зоны (\textbf{?@fig-5}).

\includegraphics[width=0.7\linewidth,height=\textheight,keepaspectratio]{image/5.png}
\href{https://github.com/ArsenioPacavira99/study_2025-2026_net-os-admin-12/tree/master/labs/lab02/report/image/5.png}{рис.
5}

Рассмотрено содержимое файлов \texttt{/var/named/named.ca},
\texttt{/var/named/named.localhost} и
\texttt{/var/named/named.loopback}: первый содержит список корневых
DNS-серверов (NS, A, AAAA-записи), второй и третий определяют локальные
зоны localhost и обратного разрешения 127.0.0.1 (\textbf{?@fig-6},
\textbf{?@fig-7}, \textbf{?@fig-8}).

\includegraphics[width=0.7\linewidth,height=\textheight,keepaspectratio]{image/6.png}
\href{https://github.com/ArsenioPacavira99/study_2025-2026_net-os-admin-12/tree/master/labs/lab02/report/image/6.png}{рис.
6}

\includegraphics[width=0.7\linewidth,height=\textheight,keepaspectratio]{image/7.png}
{[}рис.
7{]}https://github.com/ArsenioPacavira99/study\_2025-2026\_net-os-admin-12/tree/master/labs/lab02/report/image/7.png)

\includegraphics[width=0.7\linewidth,height=\textheight,keepaspectratio]{image/8.png}
\href{https://github.com/ArsenioPacavira99/study_2025-2026_net-os-admin-12/tree/master/labs/lab02/report/image/8.png}{рис.
8}

DNS-сервер запущен и добавлен в автозагрузку командами
\texttt{systemctl\ start\ named} и \texttt{systemctl\ enable\ named};
сравнение результатов \texttt{dig\ www.yandex.ru} и
\texttt{dig\ @127.0.0.1\ www.yandex.ru} показало отсутствие ответа от
локального сервера до изменения настроек прослушивания
(\textbf{?@fig-9}).

\includegraphics[width=0.7\linewidth,height=\textheight,keepaspectratio]{image/9.png}
\href{https://github.com/ArsenioPacavira99/study_2025-2026_net-os-admin-12/tree/master/labs/lab02/report/image/9.png}{рис.
9}

В NetworkManager изменены параметры соединений eth0 и System eth0:
удалены автоматические DNS, включено игнорирование auto-DNS и установлен
сервер 127.0.0.1; после перезапуска NetworkManager файл
\texttt{/etc/resolv.conf} содержит nameserver 127.0.0.1
(\textbf{?@fig-10}).

\includegraphics[width=0.7\linewidth,height=\textheight,keepaspectratio]{image/10.png}
\href{https://github.com/ArsenioPacavira99/study_2025-2026_net-os-admin-12/tree/master/labs/lab02/report/image/10.png}{рис.
10}

В файл \texttt{/etc/named.conf} внесены изменения: добавлено
прослушивание на всех интерфейсах (127.0.0.1; any;) и разрешены запросы
от сети 192.168.0.0/16; также в межсетевом экране разрешена служба DNS
(\textbf{?@fig-11}).

\includegraphics[width=0.7\linewidth,height=\textheight,keepaspectratio]{image/11.png}
\href{https://github.com/ArsenioPacavira99/study_2025-2026_net-os-admin-12/tree/master/labs/lab02/report/image/11.png}{рис.
11}

С помощью команды \texttt{lsof\ \textbar{}\ grep\ UDP} подтверждено, что
процесс named прослушивает порт 53 (UDP), что свидетельствует о
корректной работе DNS-сервера и направлении запросов через узел server
(\textbf{?@fig-12}).

\includegraphics[width=0.7\linewidth,height=\textheight,keepaspectratio]{image/12.png}
\href{https://github.com/ArsenioPacavira99/study_2025-2026_net-os-admin-12/tree/master/labs/lab02/report/image/12.png}{рис.
12}

В конфигурационном файле \texttt{/etc/named.conf} в секцию
\texttt{options} добавлены директивы
\texttt{forwarders\ \{\ 127.0.0.1;\ \};} и \texttt{forward\ first;}, а
также отключены параметры \texttt{dnssec-enable} и
\texttt{dnssec-validation}, что обеспечивает перенаправление
DNS-запросов на вышестоящий сервер и отказ от проверки DNSSEC
(\textbf{?@fig-13}).

\includegraphics[width=0.7\linewidth,height=\textheight,keepaspectratio]{image/13.png}
\href{https://github.com/ArsenioPacavira99/study_2025-2026_net-os-admin-12/tree/master/labs/lab02/report/image/13.png}{рис.
13}

\section{Конфигурирование первичного
DNS-сервера}\label{ux43aux43eux43dux444ux438ux433ux443ux440ux438ux440ux43eux432ux430ux43dux438ux435-ux43fux435ux440ux432ux438ux447ux43dux43eux433ux43e-dns-ux441ux435ux440ux432ux435ux440ux430}

На первом этапе выполнено копирование шаблона зон DNS
\texttt{named.rfc1912.zones} в каталог \texttt{/etc/named} с последующим
переименованием файла в \texttt{user.net}, после чего данный файл был
подключён в конфигурации \texttt{/etc/named.conf} директивой
\texttt{include\ "/etc/named/usernet";} (\textbf{?@fig-14}).

\includegraphics[width=0.7\linewidth,height=\textheight,keepaspectratio]{image/14.png}
\href{https://github.com/ArsenioPacavira99/study_2025-2026_net-os-admin-12/tree/master/labs/lab02/report/image/14.png}{рис.
22}

Далее в файле \texttt{/etc/named/bahis.net} вместо стандартных зон
\texttt{localhost.localdomain} и \texttt{1.0.0.127.in-addr.arpa} были
определены собственные зоны: прямая зона \texttt{bahis.net} с файлом
\texttt{master/fz/bahis.net} и обратная зона
\texttt{1.168.192.in-addr.arpa} с файлом \texttt{master/rz/192.168.1}
(\textbf{?@fig-16}).

\includegraphics[width=0.7\linewidth,height=\textheight,keepaspectratio]{image/16.png}
\href{https://github.com/ArsenioPacavira99/study_2025-2026_net-os-admin-12/tree/master/labs/lab02/report/image/16.png}{рис.
16}

После этого в каталоге \texttt{/var/named} созданы подкаталоги
\texttt{master/fz} и \texttt{master/rz} для размещения файлов прямой и
обратной зон. В каталог \texttt{master/fz} скопирован шаблон
\texttt{named.localhost}, который был переименован в \texttt{user.net}
(\textbf{?@fig-17}).

\includegraphics[width=0.7\linewidth,height=\textheight,keepaspectratio]{image/17.png}
\href{https://github.com/ArsenioPacavira99/study_2025-2026_net-os-admin-12/tree/master/labs/lab02/report/image/17.png}{рис.
17}

Затем был отредактирован файл обратной зоны
\texttt{/var/named/master/rz/192.168.1}: указана директива
\texttt{\$TTL\ 1D}, задана запись SOA с сервером
\texttt{server.user.net.}, установлен серийный номер формата ГГГГММДДВВ,
добавлена A-запись с адресом \texttt{192.168.1.1}, а также PTR-записи
для соответствия IP-адреса имени \texttt{server.bahis.net.} и
\texttt{ns.bahis.net.}; директива \texttt{\$ORIGIN} установлена в
\texttt{1.168.192.in-addr.arpa.} (\textbf{?@fig-18}).

\includegraphics[width=0.7\linewidth,height=\textheight,keepaspectratio]{image/18.png}
\href{https://github.com/ArsenioPacavira99/study_2025-2026_net-os-admin-12/tree/master/labs/lab02/report/image/18.png}{рис.
18}

Далее выполнено восстановление контекстов безопасности SELinux для
каталогов \texttt{/etc} и \texttt{/var/named}, проверены
SELinux-переключатели для службы named, после чего активирован параметр
\texttt{named\_write\_master\_zones} и сервер DNS перезапущен командой
\texttt{systemctl\ restart\ named} (\textbf{?@fig-20}).

\includegraphics[width=0.7\linewidth,height=\textheight,keepaspectratio]{image/20.png}
\href{https://github.com/ArsenioPacavira99/study_2025-2026_net-os-admin-12/tree/master/labs/lab02/report/image/20.png}{рис.
20}

\section{Проверка работоспособности первичного
DNS-сервера}\label{ux43fux440ux43eux432ux435ux440ux43aux430-ux440ux430ux431ux43eux442ux43eux441ux43fux43eux441ux43eux431ux43dux43eux441ux442ux438-ux43fux435ux440ux432ux438ux447ux43dux43eux433ux43e-dns-ux441ux435ux440ux432ux435ux440ux430}

После перезапуска службы \texttt{named} выполнена проверка разрешения
имени \texttt{ns.bahis.net} с помощью утилиты \texttt{dig}. В ответе
сервера получен статус \texttt{NOERROR}, в разделе ANSWER возвращена
A-запись \texttt{192.168.1.1}, сервером указан \texttt{127.0.0.1\#53},
что подтверждает корректную работу локального DNS-сервера
(\textbf{?@fig-21}).

\includegraphics[width=0.7\linewidth,height=\textheight,keepaspectratio]{image/21.png}
\href{https://github.com/ArsenioPacavira99/study_2025-2026_net-os-admin-12/tree/master/labs/lab02/report/image/21.png}{рис.
21}

Далее выполнена передача зоны командой \texttt{host\ -l\ user.net}. В
выводе отображаются NS-запись \texttt{server.user.net.} и A-записи для
\texttt{user.net}, \texttt{ns.user.net} и \texttt{server.user.net} с
адресом \texttt{192.168.1.1}, что подтверждает корректную настройку
прямой зоны (\textbf{?@fig-22}).

\includegraphics[width=0.7\linewidth,height=\textheight,keepaspectratio]{image/22.png}
\href{https://github.com/ArsenioPacavira99/study_2025-2026_net-os-admin-12/tree/master/labs/lab02/report/image/22.png}{рис.
22}

Командой \texttt{host\ -a\ user.net} выполнен расширенный запрос типа
ANY. В разделе ANSWER присутствуют записи SOA, NS и A, что подтверждает
корректность параметров зоны и серийного номера (\textbf{?@fig-23}).

!{[}Расширенный запрос host -a для bahis.net{]}

Затем выполнена проверка прямого и обратного разрешения имени. Команда
\texttt{host\ -t\ A\ bahis.net} возвращает IP-адрес
\texttt{192.168.1.1}. Команда \texttt{host\ -t\ PTR\ 192.168.1.1}
возвращает PTR-записи \texttt{server.bahis.net.} и
\texttt{ns.bahis.net.}, что подтверждает корректную настройку обратной
зоны \texttt{1.168.192.in-addr.arpa} (\textbf{?@fig-24}).

!{[}Проверка A- и PTR-записей зоны bahis.net{]} \#\# Внесение изменений
в настройки внутреннего окружения виртуальной машины

На виртуальной машине \texttt{server} выполнен переход в каталог
\texttt{/vagrant}, после чего созданы подкаталоги
\texttt{provision/server/dns/etc/named} и
\texttt{provision/server/dns/var/named/master}. Далее скопированы файлы
\texttt{/etc/named.conf}, содержимое \texttt{/etc/named/} и файлы зон из
\texttt{/var/named/master/} в соответствующие каталоги внутри
\texttt{/vagrant/provision/server/dns/} (рис.~\ref{fig-25}).

\begin{figure}

\centering{

\includegraphics[width=0.7\linewidth,height=\textheight,keepaspectratio]{image/25.png}

}

\caption{\label{fig-25}Подготовка структуры каталогов и копирование
конфигурации DNS в /vagrant}

\end{figure}%

Затем в каталоге \texttt{/vagrant/provision/server} создан исполняемый
файл \texttt{dns.sh}, в который внесён скрипт автоматической установки и
настройки DNS-сервера: установка пакетов \texttt{bind} и
\texttt{bind-utils}, копирование конфигурационных файлов в системные
каталоги \texttt{/etc} и \texttt{/var/named}, изменение владельца на
\texttt{named:named}, восстановление SELinux-контекстов, настройка
межсетевого экрана для службы \texttt{dns}, включение параметра
\texttt{named\_write\_master\_zones}, изменение DNS-адреса соединения
\texttt{System\ eth0} на \texttt{127.0.0.1}, перезапуск
\texttt{NetworkManager}, а также включение и запуск службы
\texttt{named} (\textbf{?@fig-26}).

\includegraphics[width=0.7\linewidth,height=\textheight,keepaspectratio]{image/26.png}
\href{https://github.com/ArsenioPacavira99/study_2025-2026_net-os-admin-12/tree/master/labs/lab02/report/image/26.png}{Нажмите
здесь, чтобы открыть ссылку.}

Для автоматического выполнения данного скрипта при загрузке виртуальной
машины в файл \texttt{Vagrantfile} в разделе конфигурации сервера
добавлен provision-блок типа \texttt{shell} с путём
\texttt{provision/server/dns.sh} и параметром
\texttt{preserve\_order:\ true}, что обеспечивает последовательное
выполнение сценария при старте VM (рис.~\ref{fig-27}).

\begin{figure}

\centering{

\includegraphics[width=0.7\linewidth,height=\textheight,keepaspectratio]{image/27.png}

}

\caption{\label{fig-27}Добавление provisioning-скрипта dns.sh в
конфигурацию Vagrantfile}

\end{figure}%

\href{https://github.com/ArsenioPacavira99/study_2025-2026_net-os-admin-12/tree/master/labs/lab02/report/image/27.png}{Нажмите
здесь, чтобы открыть ссылку.} \# Выводы

В ходе работы выполнена установка и настройка кэширующего и первичного
DNS-сервера на базе BIND.

Сервер успешно сконфигурирован для работы в режиме recursive caching, а
также в режиме первичного (master) сервера для прямой зоны
\textbf{bahis.net} и обратной зоны \textbf{1.168.192.in-addr.arpa}.

Настроены A- и PTR-записи, подтверждена авторитетность сервера по флагу
\textbf{aa} в ответах \texttt{dig}. Проверка утилитами \texttt{dig} и
\texttt{host} показала корректное разрешение имён и обратное разрешение
IP-адреса 192.168.1.1.

Обеспечена работа DNS через внутреннюю виртуальную сеть: -- изменены
параметры \texttt{listen-on} и \texttt{allow-query}; -- открыт сервис
DNS в firewall; -- настроены политики SELinux; -- сервер назначен
DNS-сервером по умолчанию (127.0.0.1).

Дополнительно выполнена автоматизация конфигурации с использованием
provisioning-скрипта \texttt{dns.sh} и интеграции в
\texttt{Vagrantfile}, что обеспечивает воспроизводимость настройки при
перезапуске виртуальной машины.

\chapter{Контрольные
вопросы:}\label{ux43aux43eux43dux442ux440ux43eux43bux44cux43dux44bux435-ux432ux43eux43fux440ux43eux441ux44b}

\begin{enumerate}
\def\labelenumi{\arabic{enumi}.}
\item
  Что такое DNS? - Это система, предназначенная для преобразования
  человекочитаемых доменных имен в IP-адреса, используемые компьютерами
  для идентификации друг друга в сети.
\item
  Каково назначение кэширующего DNS-сервера? - Его задача - хранить
  результаты предыдущих DNS-запросов в памяти. Когда клиент делает
  запрос, кэширующий DNS проверяет свой кэш, и если он содержит
  соответствующую информацию, сервер возвращает ее без необходимости
  обращаться к другим DNS-серверам. Это ускоряет процесс запроса.
\item
  Чем отличается прямая DNS-зона от обратной? - Прямая зона преобразует
  доменные имена в IP-адреса, обратная зона выполняет обратное:
  преобразует IP-адреса в доменные имена.
\item
  В каких каталогах и файлах располагаются настройки DNS-сервера? Кратко
  охарактеризуйте, за что они отвечают. - В Linux-системах обычно
  используется файл /etc/named.conf для общих настроек. Зоны хранятся в
  файлах в каталоге /var/named/, например, /var/named/example.com.zone.
\item
  Что указывается в файле resolv.conf? - Содержит информацию о DNS-
  серверах, используемых системой, а также о параметрах конфигурации.
\item
  Какие типы записи описания ресурсов есть в DNS и для чего они
  используются? - A (IPv4-адрес), AAAA (IPv6-адрес), CNAME (каноническое
  имя), MX (почтовый сервер), NS (имя сервера), PTR (обратная запись),
  SOA (начальная запись зоны), TXT (текстовая информация).
\item
  Для чего используется домен in-addr.arpa? - Используется для обратного
  маппинга IP-адресов в доменные имена.
\item
  Для чего нужен демон named? - Это DNS-сервер, реализация BIND
  (Berkeley Internet Name Domain).
\item
  В чём заключаются основные функции slave-сервера и master-сервера? -
  Master-сервер хранит оригинальные записи зоны, slave-серверы получают
  копии данных от master-сервера.
\item
  Какие параметры отвечают за время обновления зоны? - refresh, retry,
  expire, и minimum.
\item
  Как обеспечить защиту зоны от скачивания и просмотра? - Это может
  включать в себя использование TSIG (Transaction SIGnatures) для
  аутентификации между серверами.
\item
  Какая запись RR применяется при создании почтовых серверов? - MX (Mail
  Exchange).
\item
  Как протестировать работу сервера доменных имён? - Используйте команды
  nslookup, dig, или host.
\item
  Как запустить, перезапустить или остановить какую-либо службу в
  системе? - systemctl start\textbar stop\textbar restart .
\item
  Как посмотреть отладочную информацию при запуске какого-либо сервиса
  или службы? - Используйте опции, такие как -d или -v при запуске
  службы.
\item
  Где храниться отладочная информация по работе системы и служб? Как её
  посмотреть? - В системных журналах, доступных через journalctl.
\item
  Как посмотреть, какие файлы использует в своей работе тот или иной
  процесс? Приведите несколько примеров. - lsof -p или fuser -v .
\item
  Приведите несколько примеров по изменению сетевого соединения при
  помощи командного интерфейса nmcli. - Примеры включают nmcli
  connection up\textbar down .
\item
  Что такое SELinux? - Это мандатный контроль доступа для ядра Linux.
\item
  Что такое контекст (метка) SELinux? - Метка, определяющая, какие
  ресурсы могут быть доступны процессу или объекту.
\item
  Как восстановить контекст SELinux после внесения изменений в
  конфигурационные файлы? - restorecon -Rv .
\item
  Как создать разрешающие правила политики SELinux из файлов журналов,
  содержащих сообщения о запрете операций? - Используйте audit2allow.
\item
  Что такое булевый переключатель в SELinux? - Это параметр, который
  включает или отключает определенные аспекты защиты SELinux.
\item
  Как посмотреть список переключателей SELinux и их состояние? -
  getsebool -a.
\item
  Как изменить значение переключателя SELinux? - setsebool -P
  \textless on\textbar off\textgreater.
\end{enumerate}

\chapter{Список
литературы}\label{ux441ux43fux438ux441ux43eux43a-ux43bux438ux442ux435ux440ux430ux442ux443ux440ux44b}

\begin{enumerate}
\def\labelenumi{\arabic{enumi}.}
\item
  Barr D. \emph{Common DNS Operational and Configuration Errors}: RFC /
  RFC Editor. --- 02/1996. --- DOI: 10.17487/rfc1912.
\item
  McAllister M., Radvan S., Walsh D., Grift D., Paris E., Morris J.\\
  \emph{Security-Enhanced Linux. Linux с улучшенной безопасностью:
  руководство пользователя}. ---\\
  URL:
  https://docs-old.fedoraproject.org/ru-RU/Fedora/13/html/Security-Enhanced\_Linux/index.html\\
  (дата обращения: 13.09.2021).
\item
  \emph{Systemd}. --- 2015. ---\\
  URL: https://wiki.archlinux.org/index.php/Systemd\\
  (visited on 09/13/2021).
\item
  Костромин В. А. \emph{Утилита lsof --- инструмент администратора}.
  ---\\
  URL: http://rus-linux.net/kos.php?name=/papers/lsof/lsof.html\\
  (дата обращения: 13.09.2021).
\item
  Поттеринг Л. \emph{Systemd для администраторов: цикл статей}. ---
  2010. ---\\
  URL: http://wiki.opennet.ru/Systemd\\
  (дата обращения: 13.09.2021).
\item
  \emph{Сайт проекта NetworkManager}. ---\\
  URL: https://wiki.gnome.org/Projects/NetworkManager\\
  (visited on 09/13/2021).
\item
  \emph{Сайт проекта nmcli}. ---\\
  URL: https://developer.gnome.org/NetworkManager/stable/nmcli.html\\
  (visited on 09/13/2021).
\end{enumerate}


\printbibliography



\end{document}
